\chapter{Obstetric and Gynecologic Conditions}
\label{ch:obstetric}
Obstetric and gynecologic conditions encompass a wide spectrum of disorders affecting the female reproductive system across the lifespan, from menarche through pregnancy, the puerperium, and menopause. Pregnancy-related hypertensive and hemorrhagic disorders remain leading causes of maternal morbidity and mortality worldwide. Gynecologic infections, menstrual disorders, and malignancies require careful distinction from normal physiological variants.

%-------------------------------------------------------------
\section{Gynecologic Cancers (Overview)}
%-------------------------------------------------------------

%-------------------------------------------------------------

%-------------------------------------------------------------

%-------------------------------------------------------------

%-------------------------------------------------------------

%-------------------------------------------------------------

%-------------------------------------------------------------

%-------------------------------------------------------------

\label{sec:Gynecologic Cancers (Overview)}
%-------------------------------------------------------------%-------------------------------------------------------------

\subsection{Cervical Cancer}
\label{sec:obstetric:Cervical Cancer}
Cervical cancer screening with Pap smear and HPV co-testing remains the cornerstone of prevention. HPV vaccination (targeting ages 11--12, with catch-up through age 26) has the potential to nearly eliminate cervical cancer globally. Management of abnormal screening results follows the ASCCP risk-based management consensus guidelines.

\subsection{Endometrial / Uterine Cancer}
\label{sec:Endometrial Cancer}
Endometrial cancer is the most common gynecologic malignancy in developed countries. Risk factors include unopposed estrogen exposure, Lynch syndrome, advanced age, and nulliparity. Two broad histologic types exist: type I (endometrioid adenocarcinoma, 80\%, estrogen-driven, good prognosis) and type II (serous carcinoma, clear cell carcinoma, not estrogen-driven, poor prognosis). Clinical presentation is postmenopausal bleeding (90\% present with bleeding). Diagnosis is by endometrial biopsy. Treatment for early-stage disease is total hysterectomy with bilateral salpingo-oophorectomy; progestin therapy may be used for fertility-sparing management. Prognosis is generally good: 5-year survival is 95\% for stage I disease.

\subsection{Gestational Trophoblastic Disease}
\label{sec:Gestational Trophoblastic Disease}
Gestational trophoblastic disease (GTD) encompasses a spectrum of disorders arising from abnormal trophoblast proliferation, ranging from benign hydatidiform mole to malignant choriocarcinoma. Hydatidiform mole is classified as complete (46,XX, entirely paternal genome, risk of malignant transformation 15--20\%) or partial (69,XXY, triploid, risk 1--5\%). Incidence is 1 in 1000 pregnancies in Western countries. Clinical features include vaginal bleeding, exaggerated uterine size for dates, hyperemesis gravidarum, and elevated $\beta$-hCG. Ultrasonography shows a characteristic ``snowstorm'' pattern in complete mole. Treatment is suction dilation and curettage with serial $\beta$-hCG monitoring. Gestational trophoblastic neoplasia (GTN) is diagnosed when $\beta$-hCG plateaus or rises after molar evacuation; low-risk GTN is treated with single-agent methotrexate or actinomycin D; high-risk GTN requires multi-agent chemotherapy. Prognosis is excellent, with cure rates exceeding 90\% for low-risk and 80--85\% for high-risk disease.

\subsection{Ovarian Cancer}
\label{sec:obstetric:Ovarian Cancer}
Risk-reducing bilateral salpingo-oophorectomy (RRSO) is recommended for women at high risk (BRCA1 carriers by age 35--40, BRCA2 carriers by age 40--45). Oral contraceptive use reduces ovarian cancer risk by 50\% after five years of use. Screening with CA-125 and transvaginal ultrasound is not recommended for average-risk women but may be considered for high-risk women beginning age 30--35.

\subsection{Vaginal Cancer}
\label{sec:obstetric:Vaginal Cancer}
Clear cell adenocarcinoma of the vagina is a rare variant associated with in utero diethylstilbestrol (DES) exposure, affecting women born between 1940 and 1970.

\subsection{Vulvar Cancer}
\label{sec:obstetric:Vulvar Cancer}
Early recognition and biopsy of suspicious vulvar lesions (ulceration, pigmentation changes, persistent pruritus or pain) are critical for diagnosis at an early stage.

%-------------------------------------------------------------
\section{Gynecologic Infections}
%-------------------------------------------------------------

%-------------------------------------------------------------

%-------------------------------------------------------------

%-------------------------------------------------------------

%-------------------------------------------------------------

%-------------------------------------------------------------

%-------------------------------------------------------------

%-------------------------------------------------------------

\label{sec:Gynecologic Infections}
%-------------------------------------------------------------%-------------------------------------------------------------

\subsection{Bacterial Vaginosis}
\label{sec:Bacterial Vaginosis}
Bacterial vaginosis (BV) is a gynecologic infection and the most common cause of vaginal discharge in reproductive-age women, affecting 20--50\% depending on population. It is a polymicrobial dysbiosis characterized by a reduction in hydrogen peroxide-producing \textit{Lactobacillus} species and overgrowth of anaerobic bacteria. Clinical features include a thin, homogenous, gray-white vaginal discharge with a fishy odor (especially after intercourse or menses), but 50\% of women are asymptomatic. Diagnosis is by Amsel criteria ($\ge$3 of: thin white discharge, pH >4.5, positive whiff test, clue cells on saline wet mount) or Nugent score. BV in pregnancy is associated with preterm labor, PROM, and postpartum endometritis. Treatment is metronidazole (500 mg BID for 7 days or 2 g single dose) or clindamycin cream. Recurrence is common.

\subsection{Cervicitis (Chlamydia, Gonorrhea)}
\label{sec:Cervicitis}
Cervicitis is inflammation of the uterine cervix, most commonly caused by \textit{Chlamydia trachomatis} and \textit{Neisseria gonorrhoeae}. Clinical features include mucopurulent endocervical discharge, cervical friability, and intermenstrual or postcoital bleeding; many infections are asymptomatic. Untreated cervicitis can ascend to cause PID, tubal factor infertility, ectopic pregnancy, and neonatal infection. Screening with NAAT is recommended for sexually active women under 25. Treatment for chlamydia: azithromycin 1 g single dose or doxycycline 100 mg BID for 7 days; treatment for gonorrhea: ceftriaxone 500 mg IM single dose. Partner treatment and test-of-cure are indicated. Prognosis is favorable with treatment.

\subsection{Genital Herpes (HSV)}
\label{sec:Genital Herpes}
Genital herpes is caused by herpes simplex virus type 2 (HSV-2, most cases) and type 1 (HSV-1, increasing proportion). In the obstetric context, the primary concern is neonatal herpes transmission during vaginal delivery in women with active primary or recurrent genital lesions. The risk of neonatal transmission is 25--60\% with primary infection during the third trimester and less than 2\% with recurrent infection. Management at onset of labor: if active lesions or prodromal symptoms are present, cesarean delivery is recommended. For women with a history of recurrent genital herpes, suppressive acyclovir or valacyclovir is given from 36 weeks to delivery.

\subsection{Genital Warts (HPV)}
\label{sec:Genital Warts}
Genital warts are caused by low-risk human papillomavirus types 6 and 11. In the obstetric context, genital warts may increase in size and number during pregnancy. Treatment options in pregnancy include cryotherapy, trichloroacetic acid (TCA), or surgical removal; podophyllin, podofilox, and imiquimod are contraindicated. Cesarean delivery is reserved for cases where warts cause mechanical obstruction of the birth canal.

\subsection{Pelvic Inflammatory Disease}
\label{sec:obstetric:Pelvic Inflammatory Disease}
Pelvic inflammatory disease (PID) is a gynecologic infection recognized as the most common preventable cause of infertility and ectopic pregnancy worldwide. It is primarily sexually transmitted, with \textit{Chlamydia trachomatis} and \textit{Neisseria gonorrhoeae} as the predominant pathogens. The CDC-recommended treatment regimen includes ceftriaxone 500 mg IM plus doxycycline 100 mg BID for 14 days, with metronidazole for anaerobic coverage. All sexual partners should be treated.

\subsection{Trichomoniasis}
\label{sec:Trichomoniasis}
Trichomoniasis is a sexually transmitted infection caused by the protozoan \textit{Trichomonas vaginalis}, affecting approximately 170 million people worldwide. Clinical features in women include a diffuse, yellow-green, frothy, purulent vaginal discharge with a pungent odor, vulvar pruritus, dysuria, and dyspareunia; most men are asymptomatic. Trichomoniasis is associated with adverse pregnancy outcomes and increased HIV acquisition. Diagnosis is by saline wet mount showing motile trichomonads (50--70\% sensitivity), culture, or NAAT. Treatment is metronidazole 2 g single dose or tinidazole 2 g single dose, with partner treatment essential. Prognosis is favorable with treatment.

\subsection{Vulvovaginal Candidiasis}
\label{sec:Vulvovaginal Candidiasis}
Vulvovaginal candidiasis (VVC) is a fungal infection caused by \textit{Candida} species, with \textit{Candida albicans} responsible for 80--90\% of cases. Clinical features include intense vulvar pruritus, burning, vaginal soreness, dyspareunia, thick white ``cottage cheese'' discharge, vulvar erythema, fissures, and edema. Diagnosis is by saline and 10\% KOH wet mount showing pseudohyphae and budding yeast, or by culture. Treatment of uncomplicated VVC: single oral dose of fluconazole 150 mg or topical azoles. Complicated VVC requires extended therapy. Recurrent VVC ($\ge$4 episodes/year) requires maintenance therapy with fluconazole 150 mg weekly for 6 months. Prognosis is favorable with treatment.

%-------------------------------------------------------------
\section{Hemorrhagic Disorders of Pregnancy}
%-------------------------------------------------------------

%-------------------------------------------------------------

%-------------------------------------------------------------

%-------------------------------------------------------------

%-------------------------------------------------------------

%-------------------------------------------------------------

%-------------------------------------------------------------

%-------------------------------------------------------------

\label{sec:Hemorrhagic Disorders of Pregnancy}
%-------------------------------------------------------------%-------------------------------------------------------------

\subsection{Ectopic Pregnancy}
\label{sec:Ectopic Pregnancy}
Ectopic pregnancy occurs when a fertilized ovum implants outside the uterine cavity, most commonly in the fallopian tube (95\%), but also in the cervix, ovary, abdominal cavity, or previous cesarean scar. The incidence is 1--2\% of all pregnancies. The condition results from impaired tubal function or anatomy preventing normal embryo transport. Clinical features include amenorrhea, vaginal bleeding, and unilateral lower abdominal pain; ruptured ectopic pregnancy presents with sudden severe pain, syncope, hemorrhagic shock, and peritonitis. Diagnosis is by serum $\beta$-hCG and transvaginal ultrasound (empty uterine cavity, adnexal mass, free fluid); a discriminatory zone of $\beta$-hCG >2000 IU/L with no intrauterine pregnancy suggests ectopic pregnancy. Management depends on stability: stable patients with small ectopic may be treated with methotrexate; unstable patients require laparoscopic salpingectomy or salpingotomy. Prognosis is favorable with early diagnosis; ruptured ectopic pregnancy is life-threatening.

\subsection{Placenta Previa}
\label{sec:Placenta Previa}
Placenta previa is the implantation of the placenta over or near the internal cervical os, occurring in approximately 1 in 200 pregnancies at term. The condition results from abnormal trophoblast implantation in the lower uterine segment. Clinical features include painless vaginal bleeding in the second or third trimester, typically after 28 weeks gestation. Diagnosis is by transvaginal ultrasound; placenta previa is classified as complete, partial, marginal, or low-lying. Management for asymptomatic women is conservative: pelvic rest, avoidance of intercourse and vaginal exams, and elective cesarean delivery at 36--37 weeks. Massive hemorrhage may require emergency cesarean delivery, transfusion, and potentially hysterectomy. Prognosis is favorable with planned cesarean delivery.

\subsection{Placental Abruption}
\label{sec:Placental Abruption}
Placental abruption (abruptio placentae) is the premature separation of the normally implanted placenta from the uterine wall before delivery, occurring in approximately 1\% of pregnancies. The condition results from rupture of decidual spiral arteries, causing retroplacental hematoma formation and further separation. Clinical features include vaginal bleeding (80\%, though 20\% have concealed abruption), sudden sustained abdominal pain, uterine hypertonicity (``wooden'' uterus), fetal distress, and signs of hypovolemic shock that may appear out of proportion to visible blood loss. Diagnosis is primarily clinical; ultrasound has limited sensitivity (50\%). Management depends on severity and gestational age: mild abruption at term is managed by delivery; severe abruption requires immediate cesarean delivery with massive transfusion protocol. Prognosis depends on severity; complications include DIC, hemorrhagic shock, acute kidney injury, and fetal death (10--20\%).

\subsection{Postpartum Hemorrhage}
\label{sec:Postpartum Hemorrhage}
Postpartum hemorrhage (PPH) is defined as blood loss >500 mL after vaginal delivery or >1000 mL after cesarean delivery. It is the leading cause of maternal mortality worldwide, accounting for 25\% of maternal deaths. Causes are categorized by the ``four Ts'': Tone (uterine atony---70--80\% of cases), Trauma (lacerations, episiotomy, uterine rupture), Tissue (retained placenta, placenta accreta), and Thrombin (coagulopathy). Clinical features include excessive bleeding after delivery with signs of hemodynamic compromise. Diagnosis is clinical. Management follows a systematic protocol: step 1---medical management (uterine massage, oxytocin, methergine, misoprostol, carboprost); step 2---mechanical measures (uterine balloon tamponade, compression sutures); step 3---surgical intervention (laparotomy, uterine artery ligation, hysterectomy). Concurrent management includes aggressive fluid resuscitation, transfusion of blood products, and monitoring for DIC. Prognosis depends on prompt recognition and treatment.

\subsection{Uterine Rupture}
\label{sec:Uterine Rupture}
Uterine rupture is a complete disruption of all layers of the uterine wall, occurring in 0.5--1\% of women with a prior cesarean delivery attempting vaginal birth after cesarean (VBAC). It is a catastrophic obstetric emergency. The condition results from disruption of a uterine scar or traumatic delivery. Clinical features include acute onset fetal distress (fetal bradycardia most common sign---70\%), cessation of uterine contractions, loss of station of the presenting part, severe abdominal pain, and maternal hypovolemic shock. Diagnosis is often made at laparotomy. Management is immediate emergency laparotomy for delivery of the fetus and repair of the uterus; hysterectomy may be required for extensive rupture. Prognosis depends on prompt intervention; fetal mortality is 6--10\% with timely treatment.

\subsection{Vasa Previa}
\label{sec:Vasa Previa}
Vasa previa is a rare hemorrhagic disorder of pregnancy (1 in 2500 pregnancies) in which fetal blood vessels from the umbilical cord traverse the fetal membranes overlying the internal cervical os, unsupported by the placenta or umbilical cord. Fetal mortality without prenatal diagnosis is 50--60\%. Clinical features include painless vaginal bleeding at amniotomy or with spontaneous rupture of membranes, accompanied by acute fetal distress due to fetal exsanguination. Prenatal diagnosis by color Doppler ultrasound allows planned cesarean delivery at 34--35 weeks to prevent catastrophic hemorrhage. Prognosis is excellent with prenatal diagnosis.

%-------------------------------------------------------------
\section{Hypertensive Disorders of Pregnancy}
%-------------------------------------------------------------

%-------------------------------------------------------------

%-------------------------------------------------------------

%-------------------------------------------------------------

%-------------------------------------------------------------

%-------------------------------------------------------------

%-------------------------------------------------------------

%-------------------------------------------------------------

\label{sec:Hypertensive Disorders of Pregnancy}
%-------------------------------------------------------------%-------------------------------------------------------------

\subsection{Chronic Hypertension in Pregnancy}
\label{sec:Chronic Hypertension in Pregnancy}
Chronic hypertension in pregnancy is defined as hypertension predating pregnancy or diagnosed before 20 weeks of gestation. It affects 1--5\% of pregnancies. Women with chronic hypertension are at increased risk of superimposed preeclampsia (25\%), placental abruption, preterm birth, and fetal growth restriction. Management includes antihypertensives (labetalol, nifedipine, methyldopa) for BP $\ge$160/100 mmHg; ACE inhibitors and ARBs are contraindicated in pregnancy due to fetotoxicity. Prognosis is generally favorable with appropriate management.

\subsection{Eclampsia}
\label{sec:Eclampsia}
Eclampsia is a hypertensive disorder of pregnancy defined as the occurrence of generalized tonic-clonic seizures in a woman with preeclampsia that cannot be attributed to other causes. It occurs in 1 in 2000 deliveries in developed countries and more frequently in resource-limited settings. Seizures may occur antepartum (38\%), intrapartum (18\%), or postpartum (44\%). Clinical features include generalized tonic-clonic seizures in the setting of preeclampsia. Diagnosis is clinical. Management is a medical emergency: magnesium sulfate (4--6 g IV loading dose followed by 1--2 g/hour infusion) is the first-line anticonvulsant, superior to diazepam or phenytoin; airway protection, seizure precautions, and antihypertensive therapy are essential. Definitive treatment is delivery after maternal stabilization. Prognosis is favorable with prompt treatment; magnesium toxicity is treated with calcium gluconate.

\subsection{Gestational Hypertension}
\label{sec:Gestational Hypertension}
Gestational hypertension is a hypertensive disorder of pregnancy defined as new-onset hypertension (systolic BP $\ge$140 mmHg or diastolic BP $\ge$90 mmHg) after 20 weeks of gestation in the absence of proteinuria or other features of preeclampsia. It affects 6--8\% of pregnancies. The condition results from pregnancy-induced vasospasm and endothelial dysfunction. Clinical features include elevated blood pressure after 20 weeks without proteinuria. Diagnosis requires new-onset hypertension after 20 weeks in the absence of proteinuria or end-organ dysfunction. Management includes close monitoring of blood pressure and fetal surveillance; antihypertensive therapy is initiated for persistent BP $\ge$160/105 mmHg using labetalol, nifedipine, or methyldopa. Prognosis is favorable; blood pressure typically normalizes by 12 weeks postpartum, though approximately 25\% of women progress to preeclampsia.

\subsection{HELLP Syndrome}
\label{sec:HELLP Syndrome}
HELLP syndrome is a serious complication of severe preeclampsia characterized by Hemolysis, Elevated Liver enzymes, and Low Platelets. It occurs in 0.5--0.9\% of pregnancies and in 10--20\% of women with severe preeclampsia. The condition results from microvascular endothelial damage leading to platelet activation, microangiopathic hemolytic anemia, and hepatocellular necrosis. Clinical features include right upper quadrant or epigastric pain (65\%), nausea and vomiting (50\%), and headache (30\%). The Tennessee classification requires all three criteria: hemolysis (LDH >600 U/L, schistocytes on smear, bilirubin >1.2 mg/dL), elevated liver enzymes (AST $\ge$70 U/L), and low platelets (<100,000/$\mu$L). Management is delivery, which usually results in disease resolution; corticosteroids may improve platelet count temporarily; plasma exchange may be considered for persistent postpartum HELLP. Prognosis is favorable with delivery; differential diagnosis includes acute fatty liver of pregnancy, TTP, and HUS.

\subsection{Preeclampsia}
\label{sec:Preeclampsia}
Preeclampsia is a multisystem disorder of pregnancy characterized by hypertension and proteinuria developing after 20 weeks, affecting 2--8\% of pregnancies worldwide. It is a leading cause of maternal death, preterm birth, and intrauterine growth restriction. The condition results from abnormal placental implantation with inadequate spiral artery remodeling, placental ischemia, systemic endothelial dysfunction, and an imbalance of angiogenic factors (elevated soluble Flt-1, reduced PlGF). Clinical features include hypertension and proteinuria, with severe features including BP $\ge$160/110 mmHg, thrombocytopenia, elevated liver enzymes, serum creatinine >1.1 mg/dL, and persistent cerebral symptoms. Diagnosis requires BP $\ge$140/90 mmHg on two occasions four hours apart plus proteinuria $\ge$300 mg/24 hours; in the absence of proteinuria, new-onset hypertension with thrombocytopenia, renal insufficiency, impaired liver function, pulmonary edema, or cerebral symptoms may be diagnostic. Management of term preeclampsia is delivery; preterm preeclampsia is managed expectantly with antihypertensives, magnesium sulfate for seizure prophylaxis, and corticosteroids for fetal lung maturity; low-dose aspirin initiated before 16 weeks reduces risk in high-risk women. Prognosis is generally favorable with appropriate management.

%-------------------------------------------------------------
\section{Menstrual and Hormonal Disorders}
%-------------------------------------------------------------

%-------------------------------------------------------------

%-------------------------------------------------------------

%-------------------------------------------------------------

%-------------------------------------------------------------

%-------------------------------------------------------------

%-------------------------------------------------------------

%-------------------------------------------------------------

\label{sec:Menstrual and Hormonal Disorders}
%-------------------------------------------------------------%-------------------------------------------------------------

\subsection{Abnormal Uterine Bleeding}
\label{sec:Abnormal Uterine Bleeding}
Abnormal uterine bleeding (AUB) is any deviation from normal menstrual volume, duration, or regularity. The PALM-COEIN classification system provides standardized terminology: structural causes include Polyp, Adenomyosis, Leiomyoma, Malignancy and hyperplasia; non-structural causes include Coagulopathy, Ovulatory dysfunction, Endometrial, Iatrogenic, and Not yet classified. Clinical features include abnormal menstrual bleeding pattern. Diagnosis includes complete blood count, serum hCG to exclude pregnancy, TSH, prolactin, and pelvic ultrasound; hysteroscopy and endometrial biopsy are indicated for persistent AUB. Treatment depends on the cause: hormonal therapy, tranexamic acid, NSAIDs, or surgical management. Prognosis depends on the underlying cause.

\subsection{Amenorrhea (Primary, Secondary)}
\label{sec:Amenorrhea}
Amenorrhea is the absence of menstruation. Primary amenorrhea is the absence of menarche by age 15 with normal secondary sexual development, or by age 13 with no secondary sexual characteristics. Causes include hypothalamic/pituitary disorders, gonadal disorders (Turner syndrome, Swyer syndrome, premature ovarian insufficiency), and outflow tract obstruction (imperforate hymen, M\"ullerian agenesis). Secondary amenorrhea is the absence of menstruation for three months or more in a woman with previously normal cycles; the most common cause is pregnancy, with other causes including functional hypothalamic amenorrhea, PCOS, hyperprolactinemia, premature ovarian insufficiency, and Asherman syndrome. Diagnosis includes $\beta$-hCG, prolactin, TSH, FSH, and estrogen levels; progestin challenge test and MRI for pituitary lesions or outflow tract anomalies. Management depends on the underlying cause.

\subsection{Dysmenorrhea}
\label{sec:Dysmenorrhea}
Dysmenorrhea is painful menstrual cramps, the most common gynecologic complaint, affecting 50--80\% of reproductive-age women. Primary dysmenorrhea occurs without organic pelvic pathology, typically beginning within a few years of menarche, caused by excessive endometrial prostaglandin production leading to myometrial hypercontractility and uterine ischemia. Clinical features include cramping lower abdominal pain radiating to the back and thighs, beginning just before or at the onset of menstruation and lasting 1--3 days. Secondary dysmenorrhea is due to underlying pathology (endometriosis, adenomyosis, uterine fibroids, pelvic inflammatory disease). Diagnosis is clinical; secondary dysmenorrhea requires pelvic ultrasound to identify underlying causes. Treatment of primary dysmenorrhea includes NSAIDs and combined oral contraceptives. Prognosis is favorable.

\subsection{Endometriosis}
\label{sec:obstetric:Endometriosis}
Endometriosis is the most common cause of chronic pelvic pain and a leading cause of infertility. Beyond first-line therapies (NSAIDs, combined oral contraceptives), surgical excision of deeply infiltrating endometriosis may be necessary for pain refractory to medical management. GnRH agonists or antagonists with add-back therapy are used for moderate to severe disease. Dienogest and norethindrone acetate are effective progestin-only options. Treatment of endometriosis-associated infertility includes laparoscopic excision followed by IVF.

\subsection{Menopause / Perimenopause}
\label{sec:Menopause}
Menopause is the permanent cessation of menstruation for 12 consecutive months, occurring at a median age of 51 years. Perimenopause typically begins 4--8 years before menopause and is characterized by variable cycle length, vasomotor symptoms (hot flashes, night sweats), sleep disruption, and mood changes. The pathophysiology involves ovarian follicular depletion with declining inhibin B and anti-M\"ullerian hormone, followed by rising FSH and declining estradiol. Long-term consequences of estrogen deficiency include genitourinary syndrome of menopause, bone loss, and increased cardiovascular risk. Management: menopausal hormone therapy (MHT) with estrogen alone or estrogen plus progestin is the most effective treatment for vasomotor symptoms; non-hormonal options include SSRIs/SNRIs, gabapentin, and fezolinetant. Prognosis is favorable; the benefits and risks of MHT are related to age at initiation, formulation, and duration.

\subsection{Polycystic Ovary Syndrome}
\label{sec:obstetric:Polycystic Ovary Syndrome}
Polycystic ovary syndrome (PCOS) is the most common cause of anovulatory infertility and oligomenorrhea. Key management considerations include endometrial protection from unopposed estrogen: annual ultrasound assessment of endometrial thickness and biopsy if indicated, with progestins or combined oral contraceptives for cycle regulation. Long-term health risks include a 3--7-fold increased risk of endometrial cancer, 50--70\% lifetime risk of metabolic syndrome, and increased cardiovascular risk.

\subsection{Premature Ovarian Insufficiency}
\label{sec:Premature Ovarian Insufficiency}
Premature ovarian insufficiency (POI) is the loss of ovarian function before age 40, characterized by amenorrhea, elevated FSH (>40 IU/L), and low estradiol. It affects 1\% of women under 40. Causes are idiopathic (70--90\%), genetic, autoimmune, and iatrogenic. Clinical features include oligomenorrhea or amenorrhea, infertility, vasomotor symptoms, vaginal dryness, and psychological distress. POI has significant long-term health consequences due to prolonged estrogen deficiency: osteoporosis, cardiovascular disease, and cognitive decline. Management includes MHT (estrogen plus progestin continued until the average age of menopause), calcium and vitamin D supplementation, and fertility counseling. Prognosis is favorable with hormone replacement.

\subsection{Premenstrual Syndrome / Premenstrual Dysphoric Disorder}
\label{sec:PMS PMDD}
Premenstrual syndrome (PMS) is a cyclical symptom complex occurring in the luteal phase of the menstrual cycle, affecting 20--30\% of reproductive-age women. Symptoms include emotional (irritability, anxiety, depression, mood swings), physical (breast tenderness, bloating, fatigue, headaches), and behavioral changes that resolve within days of menstruation. Premenstrual dysphoric disorder (PMDD) is a severe form affecting 3--8\%, characterized by marked affective lability, irritability, depressed mood, and anxiety, causing significant functional impairment. The pathophysiology involves abnormal central nervous system response to normal cyclical fluctuations of estrogen and progesterone. Diagnosis uses DSM-5 criteria requiring prospective daily symptom charting across two cycles. First-line treatment for PMDD includes SSRIs given continuously or during the luteal phase only; other options include combined oral contraceptives and GnRH agonists for severe refractory cases. Prognosis is favorable with treatment.

%-------------------------------------------------------------
\section{Pregnancy Loss}
%-------------------------------------------------------------

%-------------------------------------------------------------

%-------------------------------------------------------------

%-------------------------------------------------------------

Pregnancy loss encompasses spontaneous abortion (miscarriage), ectopic pregnancy, recurrent pregnancy loss, and stillbirth occurring prior to viability. Etiologies include chromosomal abnormalities, uterine structural anomalies, maternal endocrine dysfunction, thrombophilias, and systemic maternal illness. Prompt evaluation requires ultrasound verification, serial beta-hCG monitoring, surgical or medical management, and compassionate psychological support.

%-------------------------------------------------------------

%-------------------------------------------------------------

%-------------------------------------------------------------

%-------------------------------------------------------------

\label{sec:Pregnancy Loss}
%-------------------------------------------------------------%-------------------------------------------------------------

\subsection{Recurrent Pregnancy Loss}
\label{sec:Recurrent Pregnancy Loss}
Recurrent pregnancy loss (RPL) is traditionally defined as two or more consecutive pregnancy losses before 20 weeks, affecting 1--2\% of couples. Evaluation is recommended after two losses and includes parental karyotyping, uterine evaluation, antiphospholipid antibody testing, and thyroid function tests. The cause remains unexplained in 50\%. Management addresses identified causes: surgical correction of uterine septum, anticoagulation for antiphospholipid syndrome, and thyroid hormone replacement for hypothyroidism; for unexplained RPL, progesterone supplementation and supportive care are associated with improved outcomes. Prognosis is variable.

\subsection{Septic Abortion}
\label{sec:Septic Abortion}
Septic abortion is a miscarriage or induced abortion complicated by infection of the uterine cavity and its contents, typically following unsafe abortion procedures. It remains a significant cause of maternal mortality in regions with restricted abortion access. Causative organisms include a mix of aerobic and anaerobic bacteria from the lower genital tract. Clinical features include fever, chills, foul-smelling vaginal discharge, lower abdominal pain, uterine tenderness, and signs of sepsis. Diagnosis is clinical with blood cultures and pelvic imaging. Treatment requires IV broad-spectrum antibiotics and prompt surgical evacuation of retained products of conception. Prognosis is favorable with prompt treatment.

\subsection{Spontaneous Abortion (Miscarriage)}
\label{sec:Spontaneous Abortion}
Spontaneous abortion, or miscarriage, is the loss of a pregnancy before 20 weeks, occurring in 10--20\% of clinically recognized pregnancies. Approximately 80\% occur in the first trimester. The most common cause (50--60\%) is chromosomal abnormalities, most often autosomal trisomies. Clinical features include vaginal bleeding with or without cramping, passage of tissue, and cervical dilation. Ultrasound evaluation determines type: threatened, inevitable, incomplete, complete, missed, or anembryonic pregnancy. Management options include expectant management, medical management (misoprostol), or surgical evacuation (dilation and curettage) for hemodynamic instability, infection, or incomplete evacuation. Prognosis is generally favorable; most early miscarriages are isolated events.

\subsection{Stillbirth (Intrauterine Fetal Demise)}
\label{sec:Stillbirth}
Stillbirth is defined as fetal death at $\ge$20 weeks (or $\ge$28 weeks by WHO definition), occurring in 1 in 160 deliveries in the United States. Causes include placental abruption, umbilical cord accidents, intrauterine infection, placental insufficiency, congenital anomalies, and maternal conditions. A cause is identified in 60--80\% of cases with thorough evaluation including fetal autopsy, placental pathology, and karyotype. Management after diagnosis is induction of labor; psychological support and counseling are critical components of care.

%-------------------------------------------------------------
\section{Pregnancy-Related Conditions}
%-------------------------------------------------------------

%-------------------------------------------------------------

%-------------------------------------------------------------

%-------------------------------------------------------------

%-------------------------------------------------------------

%-------------------------------------------------------------

%-------------------------------------------------------------

%-------------------------------------------------------------

\label{sec:Pregnancy-Related Conditions}
%-------------------------------------------------------------%-------------------------------------------------------------

\subsection{Gestational Diabetes Mellitus (Obstetric Focus)}
\label{sec:obstetric:Gestational Diabetes Mellitus}
Gestational diabetes mellitus (GDM) is a pregnancy-related condition requiring coordinated obstetric care given its implications for pregnancy outcomes. Universal screening is performed at 24--28 weeks using a 75-g oral glucose tolerance test (OGTT). Poor glycemic control in GDM is associated with an increased risk of preeclampsia, polyhydramnios, preterm labor, and cesarean delivery. Fetal surveillance includes ultrasound monitoring for macrosomia and amniotic fluid index. Intrapartum management includes close monitoring of maternal glucose to maintain normoglycemia, reducing the risk of neonatal hypoglycemia. Postpartum testing (75-g OGTT at 6--12 weeks) is essential as up to 50\% of women with GDM develop type 2 diabetes within 10 years.

\subsection{Hyperemesis Gravidarum}
\label{sec:Hyperemesis Gravidarum}
Hyperemesis gravidarum (HG) is severe, persistent nausea and vomiting during pregnancy causing weight loss (>5\% of prepregnancy weight), dehydration, electrolyte abnormalities, and ketonuria. It affects 0.5--2\% of pregnancies, distinguishing it from the more common (70--80\%) mild nausea and vomiting of pregnancy. The pathophysiology involves elevated hCG, estrogen, and progesterone, with a potential genetic predisposition (mutation in GDF15 gene). Clinical features include severe nausea and vomiting leading to weight loss and dehydration. Diagnosis is clinical. Treatment includes non-pharmacologic measures (dietary modifications, ginger, acupressure), antiemetics (doxylamine/pyridoxine combination as first-line, ondansetron, metoclopramide, promethazine), IV fluid and electrolyte repletion, and thiamine supplementation to prevent Wernicke encephalopathy. Prognosis is favorable with treatment; severe cases requiring parenteral nutrition are rare.

\subsection{Intrauterine Growth Restriction}
\label{sec:Intrauterine Growth Restriction}
Intrauterine growth restriction (IUGR), also termed fetal growth restriction (FGR), is defined as an estimated fetal weight below the 10th percentile for gestational age, often subclassified as symmetric (early onset, involving all biometric parameters) and asymmetric (late onset, with head sparing). It affects 5--10\% of pregnancies. Causes include maternal factors (hypertension, preeclampsia, chronic kidney disease, malnutrition, smoking, alcohol, illicit drugs), placental factors (abnormal placentation, infarction, velamentous cord insertion), and fetal factors (aneuploidy, congenital infections, multiple gestation). Clinical features include subnormal fetal growth on serial ultrasound. Diagnosis relies on accurate gestational dating and serial ultrasound measurement of fetal biometry, estimated fetal weight, and amniotic fluid volume; umbilical artery Doppler velocimetry identifies fetuses at highest risk. Management includes antenatal surveillance and delivery when fetal maturity is achieved or when the risks of continued in utero existence outweigh neonatal risks. Prognosis depends on the underlying cause and gestational age at delivery.

\subsection{Multiple Gestation Complications}
\label{sec:Multiple Gestation Complications}
Multiple gestation (twins 3\%, higher order multiples less common) carries increased risks of preterm birth, IUGR, preeclampsia, gestational diabetes, and placenta previa. Zygosity and chorionicity determine specific complications. Monochorionic diamniotic (MCDA) twins share a single placenta with potential vascular anastomoses, predisposing to twin-twin transfusion syndrome (TTTS, 10--15\% of MCDA twins). TTTS is diagnosed by polyhydramnios (recipient twin) and oligohydramnios (donor twin), staged by the Quintero staging system (I--V). Without treatment, severe TTTS has 80--90\% fetal mortality. Treatment includes laser photocoagulation of communicating placental vessels or serial amnioreduction. Selective intrauterine growth restriction and twin reversed arterial perfusion sequence are additional complications. Monoamniotic twins are at risk for cord entanglement and require close surveillance with planned cesarean delivery at 32--34 weeks.

\subsection{Placenta Accreta Spectrum}
\label{sec:Placenta Accreta Spectrum}
Placenta accreta spectrum (PAS) is abnormal placental adherence to the myometrium due to a deficient decidual layer, classified by depth of invasion: accreta, increta, and percreta. Incidence has risen dramatically with increasing cesarean rates, now 1 in 500--1000 pregnancies. The strongest risk factor is placenta previa in the setting of prior cesarean delivery. Clinical features include difficulty with placental separation at delivery. Prenatal diagnosis by ultrasound and MRI allows planned management. Management is controversial but increasingly favors conservative approaches when possible: planned cesarean hysterectomy with the placenta left in situ, or delayed hysterectomy after uterine artery embolization. Massive transfusion protocols and multidisciplinary teams are essential. Prognosis depends on depth of invasion and preparedness.

\subsection{Polyhydramnios and Oligohydramnios}
\label{sec:Polyhydramnios and Oligohydramnios}
Polyhydramnios is an excessive volume of amniotic fluid, defined by an amniotic fluid index (AFI) >24 cm or a deepest vertical pocket >8 cm, occurring in 1--2\% of pregnancies. Causes include idiopathic (50--60\%), maternal diabetes (10--15\%), fetal anomalies impairing swallowing, and twin-twin transfusion syndrome. Oligohydramnios is deficient amniotic fluid (AFI <5 cm or deepest vertical pocket <2 cm), affecting 1--2\% of pregnancies. Causes include fetal anomalies (renal agenesis, posterior urethral valves, obstructive uropathy), preterm PROM, placental insufficiency, and post-term pregnancy. Severe oligohydramnios in the second trimester causes pulmonary hypoplasia (Potter sequence) and limb compression deformities. Management depends on etiology and gestational age.

\subsection{Premature Rupture of Membranes}
\label{sec:Premature Rupture of Membranes}
Premature rupture of membranes (PROM) is rupture of the amniotic sac before the onset of labor. PROM at term ($\ge$37 weeks) occurs in 8\% of pregnancies and is managed with induction if labor does not ensue spontaneously within 12--24 hours. Preterm PROM (PPROM, <37 weeks) occurs in 3\% of pregnancies and is associated with ascending infection, chorioamnionitis, and preterm delivery. The condition results from weakening of the fetal membranes. Clinical features include fluid leakage from the vagina. Diagnosis is based on history, sterile speculum examination (pooling of fluid, positive Nitrazine test, ferning pattern on microscopy). Management of PPROM remote from term includes antenatal corticosteroids, antibiotics to prolong latency, and close surveillance for infection; delivery is indicated for chorioamnionitis, non-reassuring fetal status, or at 34--37 weeks. Prognosis depends on gestational age at delivery.

\subsection{Preterm Labor and Preterm Birth}
\label{sec:Preterm Labor and Preterm Birth}
Preterm labor is defined as regular uterine contractions causing cervical change between 20 and 36 weeks 6 days of gestation. Preterm birth (delivery before 37 weeks) affects 10\% of pregnancies worldwide and is the leading cause of neonatal mortality and morbidity. The pathophysiology involves multiple pathways: infection/inflammation (40\% of early preterm births), uteroplacental ischemia, uterine overdistension, and immunologically mediated rejection. Clinical features include regular contractions with cervical change before 37 weeks. Diagnosis requires documented cervical change on exam; transvaginal ultrasound measurement of cervical length (<25 mm at 20--24 weeks is high risk) and fetal fibronectin testing help stratify risk. Acute management includes antenatal corticosteroids for fetal lung maturity (24--34 weeks), magnesium sulfate for fetal neuroprotection (24--32 weeks), and tocolysis to delay delivery for 48 hours to allow corticosteroid benefit. Prognosis depends on gestational age at delivery.

\subsection{Rh Incompatibility / Hemolytic Disease of the Newborn}
\label{sec:Rh Incompatibility}
Rh incompatibility occurs when an Rh-negative mother carries an Rh-positive fetus, leading to maternal sensitization and production of anti-Rh(D) IgG antibodies that cross the placenta and cause hemolytic disease of the fetus and newborn (HDFN). Sensitization occurs most commonly at delivery but also with miscarriage, ectopic pregnancy, invasive prenatal procedures, and trauma. Clinical features range from mild hemolytic anemia to severe hydrops fetalis. Diagnosis includes positive direct antiglobulin test on cord blood, elevated bilirubin, and anemia. Management is prophylactic: Rh-immune globulin (RhoGAM, 300 $\mu$g) is administered at 28 weeks and within 72 hours of delivery of an Rh-positive infant; affected fetuses with severe anemia are treated with intrauterine transfusion. Prognosis is excellent with prophylaxis; prevention has reduced HDFN incidence by 90\%.

%-------------------------------------------------------------
\section{Puerperal Complications}
%-------------------------------------------------------------

%-------------------------------------------------------------

%-------------------------------------------------------------

%-------------------------------------------------------------

%-------------------------------------------------------------

%-------------------------------------------------------------

%-------------------------------------------------------------

%-------------------------------------------------------------

\label{sec:Puerperal Complications}
%-------------------------------------------------------------%-------------------------------------------------------------

\subsection{Bartholin Cyst}
\label{sec:Bartholin Cyst}
Bartholin cyst (and Bartholin abscess) results from fluid accumulation secondary to obstruction of the main duct of the greater vestibular (Bartholin) gland located at the 4 and 8 o'clock positions of the posterior vaginal introitus. The condition results from ductal occlusion caused by local inflammation, trauma, thickened mucus, or prior infection (Neisseria gonorrhoeae, Chlamydia trachomatis, Escherichia coli). Uninfected cysts present as painless, unilateral labial swellings. Secondary bacterial infection causes a Bartholin abscess, characterized by rapid onset of severe localized vulvar pain, dyspareunia, erythema, and a fluctuant mass. Diagnosis is clinical based on characteristic anatomic location; biopsy is indicated in women $>40$ years to rule out rare Bartholin gland carcinoma. Management of asymptomatic small cysts is conservative; symptomatic cysts or abscesses require incision and drainage with Word catheter placement or marsupialization to maintain ductal patency.

\subsection{Mastitis (Puerperal)}
\label{sec:Mastitis}
Puerperal mastitis is a common puerperal complication affecting 2--10\% of breastfeeding women, usually within the first six weeks postpartum. It presents with localized breast erythema, warmth, induration, and systemic symptoms (fever, myalgias). Management includes continued breastfeeding or milk expression to prevent milk stasis, supportive bras, and antibiotics (dicloxacillin or cephalexin). Breast abscess (5--10\% of cases) requires ultrasound-guided aspiration or incision and drainage. Prognosis is favorable.

\subsection{Postpartum Cardiomyopathy}
\label{sec:Postpartum Cardiomyopathy}
Peripartum cardiomyopathy (PPCM) is an idiopathic dilated cardiomyopathy presenting with left ventricular systolic dysfunction (LVEF <45\%) in the last month of pregnancy through five months postpartum, occurring in 1 in 1000--4000 pregnancies. The pathophysiology involves oxidative stress, a 16-kDa prolactin cleavage fragment (vasoinhibin) that impairs cardiac microvascular function, and inflammation. Clinical features include dyspnea, orthopnea, paroxysmal nocturnal dyspnea, peripheral edema, cough, and signs of left heart failure. Diagnosis is by echocardiogram showing left ventricular dilation and reduced ejection fraction. Treatment is standard heart failure therapy adjusted for pregnancy and lactation; bromocriptine may be used in some centers. Prognosis varies: 50--60\% recover LVEF to $\ge$50\% within 6 months.

\subsection{Postpartum Depression}
\label{sec:Postpartum Depression}
Postpartum depression (PPD) is a major depressive episode with onset within 4--6 weeks of childbirth (DSM-5 specifies peripartum onset). It is the most common complication of the puerperium, affecting 10--15\% of postpartum women. The pathophysiology involves a rapid drop in estrogen and progesterone after delivery, altered hypothalamic-pituitary-adrenal axis function, genetic predisposition, and psychosocial factors. Clinical features include persistent sadness, anhedonia, crying spells, fatigue, sleep and appetite disturbance, feelings of worthlessness, difficulty bonding with the infant, and in severe cases, thoughts of harming self or baby. Diagnosis is by clinical assessment using the Edinburgh Postnatal Depression Scale (EPDS, score $\ge$10). Treatment includes psychotherapy and antidepressant medication (SSRIs---sertraline and fluoxetine are first-line); for postpartum psychosis (rare, 1--2 per 1000 deliveries), emergency psychiatric care is required. Prognosis is favorable with treatment.

\subsection{Postpartum Endometritis}
\label{sec:Postpartum Endometritis}
Postpartum endometritis is infection of the decidua and myometrium after delivery, most commonly following cesarean delivery (5--30\%, versus 1--3\% after vaginal delivery). The condition results from ascending infection from the lower genital tract, most commonly polymicrobial with aerobes and anaerobes. Clinical features include fever (typically within 24--72 hours of delivery), tachycardia, uterine tenderness, and foul-smelling lochia. Diagnosis is clinical. Treatment is IV broad-spectrum antibiotics (clindamycin plus gentamicin is the standard regimen); if no response within 48 hours, add ampicillin for enterococcal coverage. Prognosis is favorable with antibiotics.

\subsection{Postpartum Thyroiditis}
\label{sec:Postpartum Thyroiditis}
Postpartum thyroiditis (PPT) is an autoimmune inflammation of the thyroid gland occurring within the first year after delivery, affecting 5--10\% of women. It is caused by a rebound of immune activity after the relative immunosuppression of pregnancy resolves, in susceptible women who are positive for thyroid peroxidase (TPO) antibodies. Clinical features include a transient thyrotoxic phase (4--8 weeks, 30\% of cases) followed by a hypothyroid phase (4--12 weeks, 40\%), with 25\% remaining hypothyroid permanently. Diagnosis is based on TSH and TPO antibody testing. The thyrotoxic phase is managed symptomatically with beta-blockers (not antithyroid drugs); the hypothyroid phase is treated with levothyroxine. Prognosis is favorable; recurrence in subsequent pregnancies is common.

